\documentclass[11pt]{article}
\usepackage{cvstyle}

\begin{document}
\cvheader



\mysection{Education}
\begin{tabular}{V}
  % 9/2014--09/2019
  2014--19 & {\bf PhD in Genome Sciences}
             \newline University of Washington, Seattle
             \newline Advisors: Michael J MacCoss, PhD and William S Noble, PhD\\
  % 9/2005--05/2009
  2005--09 & {\bf BS in Biochemistry and Molecular Biology}
             \newline Pennsylvania State University
             \newline Minor: Microbiology\\
\end{tabular}

\mysection{Employment and Professional Appointments}

\begin{tabular}{V}
  2020--     & {\bf Co-founder, Chief Technology Officer}, Talus Bioscience, Inc.\\
  2019--21   & {\bf Postdoctoral Researcher}, University of Pennsylvania
             \newline Advisor: Benjamin A Garcia, PhD\\
  2014--19 & {\bf Graduate Research Assistant}, University of Washington,
             Seattle
             \newline Advisors: Michael J MacCoss, PhD and William S Noble, PhD\\
  2011--14 & {\bf Research Associate I}, Broad Institute of MIT and Harvard
             \newline 2012--14 Proteomics Platform: Steven A Carr, PhD and Susan Abbatiello, PhD
             \newline 2011--12 Genomics Platform: Kristin Ardlie, PhD\\
  \hide{2006--09 & {\bf Undergraduate Research Assistant}, Pennsylvania State University
             \newline Advisors: Emine Koc, PhD and Melissa Rolls, PhD\\}
\end{tabular}

\mysection{Awards and Honors}
\begin{tabular}{V}
  2025     & {\bf Life Science Entrepreneurial Achievement Award}, Life Science Washington Institute  (LSWI) \\
  2024     & {\bf Rising Star Award}, Human Proteome Organization (HUPO) \\
  2021     & {\bf Rising Stars in Proteomics and Metabolomics: 40 under 40}, Journal of Proteome Research (JPR) \\
  2021     & {\bf Rising Stars in Analytical Chemistry}, Analytical Chemistry Diversity Colloquium \\
  2020     & {\bf Postdoctoral Career Development Award},  American Society for Mass Spectrometry\\ 
  2020--21   & {\bf NIH NCI T32 Training Program in Immunobiology of Normal and Neoplastic Lymphocytes}, 
            \newline NIH T32CA009140, Postdoctoral Trainee\\
  2017     & {\bf Graduate Research Award}, American Society for Mass Spectrometry\\
  2016--19 & {\bf Ruth L Kirschstein Individual National Research Service Award}
             \newline NIH F31AG055257, PI\\
             % \newline \textit{Impact Score: 25 (16\%)}\\
  \hide{2016     & {\bf NIH NIA T32 Genetic Approaches to Aging} \textit{declined}\\}
  \hide{2016     & {\bf NSF Graduate Research Fellowship Program} \textit{Honorable Mention}\\}
  2014--16 & {\bf NIH NHGRI T32 Interdisciplinary Training in Genome Sciences}
             \newline NIH T32HG000035, Predoctoral Trainee\\
  \hide{2011     & {\bf US Student Fulbright Critical Language Enhancement Award}, South Korea\\}
  2009--11  & {\bf US Fulbright Scholarship}, South Korea\\
\end{tabular}

\mysection{Research Support}
\begin{tabular}{V}
  2023--25 & {Andy Hill CARE Fund, Life Science Start-Up and Development Cycle 1 
  (Pino/Federation), Role: Principal Investigator}
             \newline \textit{\$298,973}. Awarded for ``Discovery of targeted small 
             molecule MYCN therapeutics for neuroblastoma'' \\
  2023--25 & {Andy Hill CARE Fund, Life Science Start-Up and Development Cycle 1 
  (Federation/Pino), Role: Co-Investigator}
             \newline \textit{\$1,995,337}. Awarded for ``Therapeutic disruption of the 
             Rhabdomyosarcoma transcriptional core regulatory circuit'' \\
  2023--25 & {NIH/NCI, 4R44CA268146-02 (Federation/Pino), Role: Co-Investigator}
             \newline \textit{\$1,000,000/year}. Awarded for ``Simultaneous pharmacological 
             profiling of oncogenic gene fusion proteins in cancer (FastTrack Phase II)'' \\
  2022--23 & {NIH/NCI, 1R42CA268146-01 (Federation/Pino), Role: Co-Investigator}
             \newline \textit{\$455,000}. Awarded for ``Simultaneous pharmacological 
             profiling of oncogenic gene fusion proteins in cancer (FastTrack Phase I)'' \\
  2022--23 & {NIH/NCATS, 1R43TR004221-01 (Federation/Pino), Role: Co-Investigator}
             \newline \textit{\$333,867}. Awarded for ``Illuminating the `dark' kinases -- effects 
             on proteome translocation and chromatin binding'' \\
  2021--22 & {NSF, 2112191 (Federation/Pino), Role: Co-Investigator}
             \newline \textit{\$255,831}. Awarded for ``SBIR Phase I: Simultaneous In situ 
             Pharmacological Profiling of Gene Regulatory Proteins (GRPs) for Cancer Therapy'' \\


\end{tabular}

%\newpage

\mysection{Peer-Reviewed Publications}
%{\small (X total, Y first author, Z corresponding author}
\pubcount

\begin{etaremune}

    \item Sachsenberg T, \textbf{Pino LK}, Brunet M, Bludau I, Kohlbacher O, Vizcaino JA, Bittremieux W. (2025) Perspectives in computational mass spectrometry: recent developments and key challenges. \textit{Bioinformatics Advances}. https://doi.org/10.1093/bioadv/vbaf301.

    \item Gilbert MB, Glastad KM, Fioriti M, Sorek M, Scarpa T, Purnell FS, Xu D, \textbf{Pino LK}, Korotkov A, Biashad A, Baeza J, Lauman R, Filippova A, Kacsoh BZ, Bonasio R, Mathis MW, Garcia BA, Seluanov A, Gorbunova V, Berger SL. (2025) Neuropeptides specify and reprogram division of labor in the leafcutter ant Atta cephalotes. \textit{Cell}. https://doi.org/10.1016/j.cell.2025.05.023.

    \item Shortt RL, \textbf{Pino LK}, Chea EE, Ramirez CR, Polasky DA, Nesvizhskii AI, Jones LM. (2024) Covalent Labeling Automated Data Analysis Platform for High Throughput in R (coADAPTr): A Proteome-Wide Data Analysis Platform for Covalent Labeling Experiments. \textit{J Am Soc Mass Spectrom}. https://doi.org/10.1021/jasms.4c00196.

    \item \textbf{Pino LK}. (2024) The future role of mass spectrometry in proteomics: Embracing new technologies and building bridges for higher impact. 
    \textit{J Mass Spectrom}. https://doi.org/10.1002/jms.5028
    
    \item Allen C, Meinl R, Paez JS, Searle BC, Just S, \textbf{Pino LK}, WE Fondrie. 
    (2023) nf-encyclopedia: A cloud-ready pipeline for chromatogram library 
    data-independent acquisition proteomics workflows. \textit{J Proteome Res}. 
    https://doi.org/10.1021/acs.jproteome.2c00613

    \item Lauman R, Kim HJ, \textbf{Pino LK}, Scacchetti A, Xie Y, Robison F, Sidoli S, 
    Bonasio R, Garcia BA. (2023) Expanding the Epitranscriptomic RNA Sequencing and 
    Modification Mapping Mass Spectrometry Toolbox with Field Asymmetric Waveform Ion 
    Mobility and Electrochemical Elution Liquid Chromatography. \textit{Anal Chem}. 
    https://doi.org/10.1021/acs.analchem.2c04114 

    \item Martinez TF, Lyons-Abbott S, Bookout AL, De Souza EV, Donaldson C, Vaughan 
    JM, Lau C, Abramov A, Baquero AF, Baquero K, Friedrich D, Huard J, Davis R, Kim B, 
    Koch T, Mercer AJ, Misquith A, Murray SA, Perry S, \textbf{Pino LK}, Sanford C, 
    Simon A, Zhang Y, Zipp G, Bizarro CV, Shokhirev MN, Whittle AJ, Searle BC, MacCoss 
    MJ, Saghatelian A, Barnes CA. (2023) Profiling mouse brown and white adipocytes to 
    identify metabolically relevant small ORFs and functional microproteins. \textit{Cell 
    Metab}. https://doi.org/10.1016/j.cmet.2022.12.004

    \item \textbf{Pino LK}, Schilling B. (2021) Proximity labeling and other novel mass 
    spectrometric approaches for spatiotemporal protein dynamics. \textit{Expert Rev 
    Proteom}. https://doi.org/10.1080/14789450.2021.1976149

    \item Wang SY, Pollina EA, Wang IH, \textbf{Pino LK}, Bushnell HL, Takashima K, 
    Fritsche C, Sabin G, Garcia BA, Greer PL, Greer EL. (2021) Role of epigenetics 
    in unicellular to multicellular transition in Dictyostelium. \textit{Genome Biol}. 
    https://doi.org/10.1186/s13059-021-02360-9

    \item \textbf{Pino LK}, Baeza J, Lauman R, Schilling B, Garcia BA. (2021) Improved 
    SILAC quantification with data independent acquisition to investigate 
    bortezomib-induced protein degradation. \textit{J Proteome Res}. 
    https://doi.org/10.1021/acs.jproteome.0c00938

    \item \textbf{Pino LK}, Rose J, O’Broin A, Shah S, Schilling B. (2020) Emerging mass
    spectrometry-based proteomics methodologies for novel biomedical applications.
    \textit{Biochem Soc Trans} https://doi.org/10.1042/BST20191091
    
    \item Chiao YA, Zhang H, Mariya Sweetwyne M, Jeremy Whitson J, Ting YS, Nathan Basisty
    N, \textbf{Pino LK}, Ellen Quarles E, Nguyen NH, Campbell MD, Zhang T, Gaffrey MJ, 
    Merrihew G, Wang L, Yue Y, Duan D, Granzier HL, Szeto HH, Qian WJ, Marcinek D, MacCoss 
    MJ, Rabinovitch P. (2020) Late-life restoration of mitochondrial function reverses 
    cardiac dysfunction in old mice. \textit{eLife} https://doi.org/10.7554/eLife.55513
    
    \item \textbf{Pino LK}, Just S, MacCoss MJ, Searle BC. (2020) Acquiring and analyzing 
    data independent acquisition proteomics experiments without spectrum libraries. 
    \textit{Mol Cell Proteomics} https://doi.org/10.1074/mcp.P119.001913
    
    \item Federation AJ, Nandakumar V, Searle BC, Stergachis A, Wang H, \textbf{Pino LK}, 
    Merrihew G, Ting YS, Howard N, Kutyavin T, MacCoss MJ, Stamatoyannopoulos JA. (2020) 
    Highly parallel quantification and compartment localization of transcription factors 
    and nuclear proteins. \textit{Cell Reports} 
    https://doi.org/10.1016/j.celrep.2020.01.096
    
    \item \textbf{Pino LK}, Searle BC, Yang HY, Hoofnagle AN, Noble WS, MacCoss MJ. 
    (2020) Matrix-matched calibration curves for assessing analytical figures of merit 
    in quantitative proteomics. \textit{J Proteome Res}
    https://doi.org/10.1021/acs.jproteome.9b00666

    \item \textbf{Pino LK}, Lin A, Bittremieux W. (2019) 2018 YPIC Challenge: A Case
    Study in Characterizing an Unknown Protein Sample. \textit{J Proteome Res}
    https://doi.org/10.1021/acs.jproteome.9b00384

    \item Searle BC, \textbf{Pino LK}, Egertson JD, Ting YS, Lawrence RL, MacLean BX,
    Villen J, MacCoss MJ. (2018) Chromatogram libraries improve peptide detection and 
    quantification by data independent acquisition mass spectrometry. \textit{Nat Comm} 
    Dec 2018. https://doi.org/10.1038/s41467-018-07454-w
    
    \item \textbf{Pino LK}, Searle BC, Huang E, Noble WS, Hoofnagle AN, MacCoss MJ. 
    (2018) Calibration using a single-point external reference material harmonizes 
    quantitative mass spectrometry proteomics data between platforms and laboratories. 
    \textit{Anal Chem} https://doi.org/10.1021/acs.analchem.8b04581
    
    \item Galitzine C, Egertson JD, Abbatiello S, Henderson CM, \textbf{Pino LK}, 
    MacCoss M, Hoofnagle AN, Vitek O. (2018) Nonlinear regression improves accuracy of 
    characterization of multiplexed mass spectrometric assays. \textit{Mol Cell Prot} 
    https://doi.org/10.1074/mcp.RA117.000322
    
    \item \textbf{Pino LK}, Searle BC, Bollinger JG, Nunn B, MacLean BX, MacCoss MJ. 
    (2017) The Skyline Ecosystem: Informatics for Quantitative Mass Spectrometry 
    Proteomics. \textit{Mass Spectrom Rev} https://doi.org/10.1002/mas.21540
    
    \item Abelin JG, Patel J, Lu X, Feeney CM, Fagbami L, Creech AL, Hu R, Lam D, Davison 
    D, \textbf{Pino L}, Qiao JW, Kuhn E, Officer A, Li J, Abbatiello S, Subramanian A, 
    Sidman R, Snyder E, Carr SA, Jaffe JD. (2016) Reduced-representation 
    Phosphosignatures Measured by Quantitative Targeted MS Capture Cellular States and 
    Enable Large-scale Comparison of Drug-induced Phenotypes. \textit{Mol Cell Proteomics} 
    https://doi.org/10.1074/mcp.M116.058354
    
\end{etaremune}

\mysection{Preprints}
\begin{etaremune}

    \item Wu A, Kohler D, Navada PP, Robbins J, Boyle G, Boshart A, Karis K, Neefjes J, Konvalinka A, Sarthy JF, \textbf{Pino LK}, Gyori BM, Vitek O. (2026) MSstatsBioNet: Integrating Statistical Analyses with Prior Knowledge Biomolecular Networks for Quantitative Proteomics and Phosphoproteomics. \textit{bioRxiv}. https://doi.org/10.64898/2026.07.09.737605.

    \item Wen B, Paez JS, Hsu C, Canzani D, Chang A, Shulman N, MacLean BX, Berg MD, Villen J, Fondrie WE, \textbf{Pino LK}, MacCoss MJ, Noble WS. (2026) Carafe2 enables high quality \textit{in silico} spectral library generation for timsTOF data-independent acquisition proteomics. \textit{bioRxiv}. https://doi.org/10.64898/2026.03.27.714846.

    \item \textbf{Pino LK*}, Canzani D, Gutierrez A, Robbins J, McEllin B, Hubbard E, Broderick E, Prymolenna A, Tatka L, Paez JS, Mercenne G, Siebenthall K, Fondrie WE, Federation AJ*. \textit{[*co-corresponding]} (2025) Native, Spatiotemporal Profiling of the Global Human Regulome. \textit{bioRxiv}. https://doi.org/10.1101/2025.06.14.659727.

\end{etaremune}

\mysection{Additional Publications}
\begin{etaremune}
  \item {\bf Pino LK}. Methods for harmonizing and calibrating quantitative mass spectrometry 
  experiments. PhD dissertation. University of Washington, Seattle. Aug, 2019. Advisors: 
  Michael J MacCoss, PhD and William S Noble, PhD
\end{etaremune}

\mysection{Invited Research Lectures}
\begin{etaremune}

% Title as supplied, not read off a slide: the deck
% (20260805_proteinsocietythailand.pptx) is 82 MB and returns no extractable text.
    \item \textbf{Pino LK}. Proteomics to drive drug discovery through direct measurement of transcriptional regulation. (Aug 2026) \textit{The Protein Society}, Thailand.

    \item \textbf{Pino LK}. Accidental Entrepreneur: Founding a Biotech Startup Company. (Jul 2026) \textit{Seattle University Summer Accelerator}, Seattle, WA.

    \item \textbf{Pino LK}. Integrative regulome profiling of DNA damage responses to discover DNA repair-directed therapeutics in aging. (Jun 2026) \textit{American Society for Mass Spectrometry 2026 Annual Conference}, San Diego, CA, selected oral presentation.

    \item \textbf{Pino LK}. Chromatogram-based chemoproteomics with the ZenoTOF 8600 system and Skyline. (Jun 2026) \textit{Sciex Breakfast Seminar, American Society for Mass Spectrometry 2026 Annual Conference}, San Diego, CA.

% VENUE? Still open whether this was in person at Stanford or virtual. Title as
% supplied; this deck (20260514_stanford-ugm.pptx) is also too large to extract.
    \item \textbf{Pino LK}. Scaling Proteomics to Unlock Transcription Factors for Drug Discovery. (May 2026) \textit{Stanford Proteomics Users Group (PUG)}, Stanford, CA.

% Same title as the Mar 2026 ACS talk below: the LSINW deck is a copy of the ACS
% deck whose title slide was never updated (it still reads "ACS 2026 Division of
% Medicinal Chemistry, 2026 March 24"). Accurate as far as the deck goes, but you
% may want to differentiate the two.
    \item \textbf{Pino LK}. Unlocking the regulo.me for drug discovery using a high-throughput proteomics platform. (Apr 2026) \textit{Life Science Innovation NorthWest, AI session}, Seattle, WA.

    \item \textbf{Pino LK}. From PhD to Platform: Founding a Biotech Startup Company. (Apr 2026) \textit{Bioscience Careers Seminar Series, University of Washington PhD Ambassadors}, Seattle, WA.

% Location inferred from the deck filename (20260402_sciex_pharma-seattle.pptx);
% confirm it was Seattle and not a webinar.
    \item \textbf{Pino LK}. SCIEX ZenoTOF 8600 system in chemoproteomics drug discovery applications. (Apr 2026) \textit{Discovery Continuum Seminar, SCIEX}, Seattle, WA.

    \item \textbf{Pino LK}. Unlocking the regulo.me for drug discovery using a high-throughput proteomics platform. (Mar 2026) \textit{American Chemical Society Division of Medicinal Chemistry, session ``Drugging Transcription Factors: From Privileged Classics to the Next Frontier''}, Atlanta, GA.

% TITLE? The last remaining placeholder. No deck for this one anywhere in the
% presentations archive, so the title has to come from memory, the event page, or
% your calendar. If it cannot be found, wrap this entry in \hide{} rather than
% printing the placeholder.
    \item \textbf{Pino LK}. Keynote address. (Mar 2026) \textit{2026 Hollomon Health Innovation Challenge}, Seattle, WA.

    \item \textbf{Pino LK}. Chromatogram-based chemoproteomics with ZenoTOF 8600 and Skyline. (Feb 2026) \textit{Sciex Lunch Seminar, US HUPO 2026 Annual Conference}, St.\ Louis, MO.

    \item \textbf{Pino LK}. Small-molecule modulation of chromatin-associated protein function in native cells. (Jan 2026) \textit{University of Washington Department of Chemistry}, Seattle, WA.

    \item \textbf{Pino LK}. Profiling the regulome with high-throughput proteomics to enable small-molecule therapeutics. (Dec 2025) \textit{Northwestern University Department of Biochemistry and Molecular Genetics}, Chicago, IL.

    \item \textbf{Pino LK}. From regulome maps to first-in-class targets through functional proteomics at scale. (Nov 2025) \textit{Bruker Lunch Seminar, 24th Human Proteome Organization World Congress}, Toronto, Canada.

    \item \textbf{Pino LK}. Spatiotemporal Dynamics of the Transcriptional Regulome Across Cell State and Drug Response Landscapes. (Nov 2025) \textit{24th Human Proteome Organization World Congress}, Toronto, Canada.

% The Bruker "Bringing ASMS 2025 to Seattle" event -- TODO.md had you confirming
% you spoke. The deck's title slide carries only the date, no venue, so the event
% name comes from the filename (20250722_bruker_scaling).
    \item \textbf{Pino LK}. Proteomics for Chromatin-Bound Protein Drug Discovery That Doesn't Fall Apart at Scale. (Jul 2025) \textit{Bruker ``Bringing ASMS 2025 to Seattle''}, Seattle, WA.

% Accelerator pitch rather than a research lecture -- hidden here; the accelerator
% itself may belong under Entrepreneurship instead.
    \hide{\item \textbf{Pino LK}. The regulome has been hidden. Talus reveals it. (Jul 2025) \textit{AWS Generative AI Accelerator}, virtual.}

% Third stop on the same Bruker ABPP roadshow (see Sep 2024 San Francisco and Nov
% 2024 San Diego) -- identical title, hidden as a repeat.
    \hide{\item \textbf{Pino LK}. Activity-based Protein Profiling to Identify Functional Effects of Small Molecules on Transcription Factors. (Mar 2025) \textit{Bruker User Meeting}, Los Angeles, CA.}

    \item \textbf{Pino LK}. Drug Discovery Screening for Transcription Factors from Low-Input, High-Throughput, Cell-Based Proteomics. (Feb 2025) \textit{US HUPO 2025 Annual Conference, Session 9 ``Down in the Dumps: Treasure from Troubleshooting''}, Philadelphia, PA.

% Evening workshop rather than a lecture -- you build lab startup budgets with the
% room. Might sit better under Teaching.
    \hide{\item \textbf{Pino LK}. The Business of Starting a Lab Business. (Feb 2025) \textit{Evening Workshop, US HUPO 2025 Annual Conference}, Philadelphia, PA.}

% Second stop on the Bruker ABPP roadshow -- identical title to Sep 2024 San
% Francisco, hidden as a repeat. The middle section was reworked to brachyury and
% chordoma, so unhide if you would rather list it as a distinct talk.
    \hide{\item \textbf{Pino LK}. Activity-based Protein Profiling to Identify Functional Effects of Small Molecules on Transcription Factors. (Nov 2024) \textit{Bruker User Meeting}, San Diego, CA.}

    \item \textbf{Pino LK}. From Lab to Launch: Navigating the Entrepreneurial Path in Start-up Biotech. (Nov 2024) \textit{Pharmaceutical Development and Preclinical Discovery, American Association of Pharmaceutical Scientists}, virtual.

% Answers RECALL.md's question for 2024: yes, the Rising Star Award came with a
% lecture. The title slide carries no date and does not name Dresden, so both come
% from the deck filename and the surrounding HUPO 2024 talks.
    \item \textbf{Pino LK}. It Takes a Village: Building a proteomics community for fun and profit. (Oct 2024) \textit{Rising Star Award lecture, 23rd Human Proteome Organization World Congress}, Dresden, Germany.

    \item \textbf{Pino LK}. Integrating AI throughout a Mass Spectrometry-driven Drug Discovery Approach. (Oct 2024) \textit{23rd Human Proteome Organization World Congress}, Dresden, Germany.

    \item \textbf{Pino LK}. Transcription Factors on Target: Leveraging Targeted Mass Spectrometry Proteomics for High-Throughput Small Molecule Screening. (Oct 2024) \textit{Thermo Lunch Symposium ``New Advancements with the Thermo Scientific Stellar MS'', 23rd Human Proteome Organization World Congress}, Dresden, Germany.

% Date conflict: the filename says Oct 20, the title slide says Oct 22. The
% workshop was pre-congress, so Oct 20 is the likelier of the two.
    \item \textbf{Pino LK}. Proteomics driven drug discovery in a startup company. (Oct 2024) \textit{Pre-congress Workshop 2: Chemical proteomics and drug discovery, 23rd Human Proteome Organization World Congress}, Dresden, Germany.

% First of three stops on the Bruker ABPP roadshow, shown as the representative one;
% San Diego (Nov 2024) and Los Angeles (Mar 2025) are hidden as repeats.
    \item \textbf{Pino LK}. Activity-based Protein Profiling to Identify Functional Effects of Small Molecules on Transcription Factors. (Sep 2024) \textit{Bruker User Meeting}, San Francisco, CA.

    \item \textbf{Pino LK}. Developing and optimizing mass spectrometry methods to improve detection of covalently modified peptides in complex lysates. (Jun 2024) \textit{American Society for Mass Spectrometry 2024 Annual Conference}, Anaheim, CA.
    
    \item \textbf{Pino LK}. Profiling oncogenic transcription factors using mass spectrometry proteomics. (Jun 2024) \textit{Bruker eXceed Symposia 2024}, Anaheim, CA.

    \item \textbf{Pino LK}. Embracing Disruption and New Opportunities in the Evolution of Proteomics Technologies. (Apr 2024) \textit{Association of Biomolecular Resource Facilities}, Minneapolis, MN.

% sPRG research-group talk rather than a conference lecture, and the date is
% contested: the filename says Apr 23 but the title slide reads "Cold Spring Harbor
% Laboratory, 2024 Aug 10", so the deck looks reused for the CSHL course. Hidden
% until you settle which occasion this was.
    \hide{\item \textbf{Pino LK}. Benchmarking and Sanity Checks for LFQ. (Apr 2024) \textit{Proteomics Standards Research Group, Association of Biomolecular Resource Facilities}, Minneapolis, MN.}

    \item \textbf{Pino LK}. From Lab to Launch: Navigating the Entrepreneurial Path in Start-up Biotech. (Mar 2024) \textit{Ohio State University Biochemistry Department}, virtual.

    \item \textbf{Pino LK}. Transcription Factor Profiling by Mass Spectrometry to Discover and Develop Cancer Therapeutics. (Mar 2024) \textit{US HUPO Parallel Session 14: Transcription Factors}, Portland, OR.

    \item \textbf{Pino LK}. Automating proteomics for scale: Chasing asymptotic efficiency. (Mar 2024) \textit{US HUPO Bruker Sponsored Lunch}, Portland, OR.

% A three-hour short course, so arguably Teaching rather than an invited lecture.
    \hide{\item \textbf{Pino LK}. Pitching Proteomics: Proteomics for Non-scientists. (Mar 2024) \textit{Short Course, US HUPO Annual Conference}, Portland, OR.}

% DATE? The CV and RECALL.md both say Jan 2024, but the deck
% (20240221_sfu-vancouver_keynote.pptx) prints "February 21, 2024" on its title
% slide. Worth settling which is right.
    \item \textbf{Pino LK}. Mass Spectrometry-Based Approaches for Chemoproteomics and
    Activity-Based Protein Profiling. (Jan 2024) \textit{Annual BCPM Network Symposium},
    Simon Fraser University, Canada.

    \item \textbf{Pino LK}. Mass Spectrometry-Based Approaches for Chemoproteomics and 
    Activity-Based Protein Profiling. (Jan 2024) \textit{Bruker eXceed Event: UC 
    Berkeley Mass Spec Seminar}, Berkeley CA.

    \item \textbf{Pino LK}. Profiling the DNA Regulome to Discover Direct Inhibitors of
    the Brachyury Transcription Factor. (Nov 2023) \textit{Ohio State University
    Biochemistry Seminar}, Columbus OH.

% Same title as the two Jan 2024 entries above -- this Asilomar talk is where it
% started.
    \item \textbf{Pino LK}. Mass spectrometry-based approaches for chemoproteomics and activity-based protein profiling. (Oct 2023) \textit{38th Asilomar Conference on Mass Spectrometry: Computational Mass Spectrometry}, Pacific Grove, CA.

% Answers RECALL.md's 2023 question about a talk tied to the CARE and NIH awards.
    \item \textbf{Pino LK}. Profiling the DNA Regulome to Discover Direct Inhibitors of the Brachyury Transcription Factor. (Sep 2023) \textit{Discovery on Target: Targeting Transcription Factors}, Boston, MA.

    \item \textbf{Pino LK}. Quantitative Accuracy at Speed with the timsTOF SCP. (Jun 2023) \textit{Bruker Users Meeting for Life Science Mass Spectrometry}, Seattle, WA.

% Answers RECALL.md's 2023 question: yes, you spoke at ASMS 2023. Section dividers
% inside the deck use a different working title, "A scalable mass spectrometry
% proteomics-based platform to screen for transcription factor small molecule
% therapeutics" -- the title slide version is used here.
    \item \textbf{Pino LK}. Simultaneous In situ Pharmacological Profiling of Transcription Factors for Cancer Therapy. (Jun 2023) \textit{American Society for Mass Spectrometry 2023 Annual Conference}, Houston, TX.

    \item \textbf{Pino LK}. Epigenetics and Entrepreneurship: Proteomics for everything from developmental biology to cancer drug discovery. (Mar 2023) \textit{Brigham Young University}, Provo, UT.

    \item \textbf{Pino LK}. Prospective system suitability and quality control design for
    large-scale and longitudinal quantitative mass spectrometry proteomics. (Mar 2022)
    \textit{US HUPO Annual Conference}, Charleston, SC.

% Same title as the US HUPO talk above, delivered again at LabKey -- hidden as a repeat.
    \hide{\item \textbf{Pino LK}. Prospective system suitability and quality control design for large-scale and longitudinal quantitative mass spectrometry proteomics. (Oct 2022) \textit{LabKey Conference}.}

% VENUE? The filename says US HUPO 2022 but the title slide names no venue and
% prints Mar 01 2022 rather than Feb 28. Hidden until the venue is confirmed.
    \hide{\item \textbf{Pino LK}. Using Data Independent Acquisition to Improve Quantification for Spatiotemporal Proteomics. (Feb 2022) \textit{US HUPO Annual Conference}, Charleston, SC.}

% Partner meeting rather than an invited lecture.
    \hide{\item \textbf{Pino LK}. Comprehensive, Quantitative Measurement of Proteins Participating in Gene Regulation. (Dec 2021) \textit{Amicus Therapeutics}.}

% VENUE? The filename says "TF conf" and the speaker notes mention a virtual
% conference, but no venue, date or organizer appears anywhere in the deck.
    \hide{\item \textbf{Pino LK}. Comprehensive, Quantitative Measurement of Proteins Participating in Gene Regulation. (Jun 2021) \textit{Transcription factor conference}, virtual.}

    \item \textbf{Pino LK}, Baeza J, Lauman RA, Garcia BA. (Jun 2020) Opportunities and 
    challenges for improving quantitative accuracy and precision in SILAC with 
    DIA-MS. \textit{ASMS Annual Conference}, virtual.
    
    \item \textbf{Pino LK}, Searle BC, Yang HY, Hoofnagle AN, Noble WS, 
    MacCoss MJ. (Jun 2019) Using an external reference material to harmonize and calibrate 
    quantitative mass spectrometry data at scale. \textit{ASMS Annual 
    Conference}, Atlanta, GA.
    
    \item \textbf{Pino LK}, Searle BC, Hoofnagle AN, Noble WS, MacCoss MJ. (Jun 2018) 
    Signal calibration for quantitative proteomics. \textit{Skyline User Group Meeting}, 
    San Diego, CA.
    
    \item \textbf{Pino LK}, Searle BC, Egertson JD, Yang HY, Noble WS, MacCoss MJ. (Jul 
    2016) Reproducible quantification of the yeast proteome using DIA mass spectrometry. 
    \textit{Cascadia Proteomics Symposium}, Seattle, WA.
    
    \item \textbf{Pino LK}, Searle BC, Egertson JD, Yang HY, Noble WS, MacCoss MJ. (Jun 
    2016) Applying lessons learned from targeted mass spectrometry to data-independent 
    acquisition (DIA) assays. \textit{Skyline User Group Meeting}, San Antonio, TX.
    
\end{etaremune}

%\newpage

\mysection{Panels, Moderation, and Judging}
\begin{etaremune}

    \item Panelist, ``Aging in the Era of AI'' happy hour. (Jul 2026) \textit{Seattle Tech Week}, Seattle, WA.

    \item Moderator, Chemoproteomics workshop panel. (Jun 2026) \textit{American Society for Mass Spectrometry 2026 Annual Conference}, San Diego, CA.

    \item Moderator, Entrepreneurship workshop panel. (Jun 2026) \textit{American Society for Mass Spectrometry 2026 Annual Conference}, San Diego, CA.

    \item Panelist, Biotech Panel. (Apr 2026) \textit{University of Washington Project Medicine and FosterMed}, Seattle, WA.

    \item Panelist, ``From Lab to Launch''. (Apr 2026) \textit{Event for ITI, Seabridge, and WRF Fellows}, Seattle, WA.

    \item Judge. (Mar 2026) \textit{2026 Hollomon Health Innovation Challenge}, Seattle, WA.

    \item Panelist, Early Career Researcher panel. (Mar 2023) \textit{Human Proteome Organization}, virtual.

% Your contribution is a single "Call to Action / Take Home Message" slide titled
% "Functional proteomics for phenotypic drug discovery".
    \item Panelist, Thermo Visionary Panel Discussion. (Jun 2022) \textit{American Society for Mass Spectrometry 2022 Annual Conference}, Minneapolis, MN.

\end{etaremune}

% Every poster hidden for now -- the whole section is wrapped, heading included,
% so the curated build does not print an empty heading.
\hide{\mysection{Posters}
\begin{etaremune}

    \item \textbf{Pino LK}, Searle BC, Hoofnagle AN, Noble WS, MacCoss MJ. Development and validation of a multiplex quantitative protein assay for human cerebrospinal fluid. \textit{ASMS Annual Conference}, Jun 2018, San Diego, CA.
    
    \item \textbf{Pino LK}, Kaeberlein M, Hoofnagle AN, Noble WS, MacCoss MJ. Lessons learned in high-dimensional analysis of quantitative DIA-MS proteome profiling. \textit{Cascadia Proteomics Symposium}, Aug 2017, Seattle, WA.
    
    \item \textbf{Pino LK}, Kaeberlein M, Noble WS, MacCoss MJ. Systematic proteome profiling by DIA-MS to characterize calorie restriction-induced life span extension in yeast. \textit{ASMS Annual Conference}, Jun 2017, Indianapolis, IN.
    
    \item \textbf{Pino LK}, Merrihew G, Noble WS, MacCoss MJ. Reproducible quantification of the yeast proteome using data independent acquisition (DIA) mass spectrometry. \textit{ASMS Annual Conference}, Jun 2016, San Antonio, TX.
    
    \item \textbf{Pino LK}, Abbatiello SE, Fagbami L, Greulich H, Kuhn E, Jaffe J, Carr SA. New Challenges in Targeted Peptide Quantitation: Phosphopeptides. \textit{ASMS Annual Conference}, Jun 2014, Baltimore, MD.
    
    \item \textbf{Pino LK}, Abbatiello SE, Belford M, Kuhn E, Yates N, Carr SA. Improving Speed and Selectivity of Targeted Peptide Quantification using FAIMS. \textit{ASMS Annual Conference}, Jun 2013, Minneapolis, MN.

  \end{etaremune}}

\mysection{Mentoring}
\begin{tabular}{V}
  \raggedright
  2025--26 & {\bf Devon Kohler}, PhD student, Northeastern University
            \newline Thesis: Statistical and causal inference methods improve accuracy and interpretability in biomolecular experiments\\
  2025     & {\bf Nina Shetty}, undergraduate, University of Washington
            \newline Buerk Center's Women's Entrepreneurial Leadership program (WE Lead)\\
  2020--22 & {\bf Gabriela Witek}, PhD student, University of Pennsylvania
            \newline Thesis: Proteomics of ALK mutation in neuroblastoma\\
  2020     & {\bf Adetola Alonge}, undergraduate, University of Florida
            \newline Project: Metaanalysis of Patient Cerebrospinal Fluid by Mass Spectrometry Proteomics Discovers Known and Novel Biomarkers of Parkinson's Disease
            \newline \textit{ePoster Award for Biochemistry and Molecular Biology, Annual Biomedical Research Conference for Minority Students (ABRCMS) (Nov 2020)} \\
  2019--24 & {\bf Michael Gilbert}, PhD student, University of Pennsylvania
            \newline Thesis: Investigating Caste Differences in \textit{Atta cephalotes}\\
\end{tabular}

% \newpage
\mysection{Entrepreneurship}

\mysubsection{Accelerator Programs}
\begin{tabular}{V}
  2025 & \textit{Founder To Leader}. Femme Fellowship. \\
  2022 & \textit{Creative Destruction Laboratories (CDL)}. Computational Health (Seattle) graduate. \\
  2022 & \textit{Project W}. Women Entrepreneurs Boot Camp. \\
  2021 & \textit{Y Combinator (YC), S21 Batch}. Group partners: Surbhi Sarna, Jared Friedman, Brad Flora. \\
  2020--24 & \textit{Life Sciences Washington Institute}. Washington Innovation Network (WIN) graduate.\\
\end{tabular}

% Wrapped whole: hiding only the rows would leave a heading over an empty table.
\hide{\mysubsection{Company Appearances}
\begin{tabular}{V}
  2026 & \textit{Seattle Tech Week Startup Showcase Fair}. Featured company table, Seattle, WA.\\
\end{tabular}}

% \newpage
\mysection{Professional Service}

\mysubsection{Editorial Boards}
\begin{tabular}{V}
  2021--23 & \textit{Journal of Proteome Research}. Editorial Advisory Board. \\
  2021 & \textit{Molecular Omics}. Guest editor, special issue covering the 2021 Annual US HUPO Conference. \\
  2021 & \textit{Molecular Omics}. Guest editor, special issue titled ``Cross-talk and Integrative Post-translational Modifications''. \\
\end{tabular}

{\bf Journal Referee: }{Nature Communications, Molecular \& Cellular Proteomics, Journal of Proteome Research, Analytical Chemistry, Scientific Reports, Proteomics, Frontiers in Molecular Neuroscience, Journal of the American Society for Mass Spectrometry, Clinical Proteomics}\\ 

\mysubsection{Committee and Research Group Membership}
\begin{tabular}{V}
  2023--24 & Member-at-Large, US Human Proteome Organization (US HUPO). \textit{Board of Directors}. \\
  2021--24 & Proteomics Standards Research Group (sPRG), Association of Biomolecular Resources Facilities (ABRF). \textit{Committee member}. \\
  2021-- & Computational Mass Spectrometry (CompMS) Community of Special Interest (COSI), Intelligent Systems for Molecular Biology Conference (ISMB). \textit{Co-chair}. \\
% Description(s) of tasks/responsibilities
        %\item {Organize and facilitate interest group meetings and workshops during the annual conference.}
  2020--22 & Early Career Research Group, US Human Proteome Organization (US HUPO). \textit{Committee member}. \\
% Description(s) of tasks/responsibilities
        %\item {Contribute to organizing committee for 2021 US HUPO conference.}
  2020--23 & Data Independent Acquisition Interest Group, American Society for Mass Spectrometry (ASMS). \textit{Co-coordinator}. \\
  2019--21 & Digital Communications Committee, American Society for Mass Spectrometry (ASMS). \textit{Committee member}. \\
% Description(s) of tasks/responsibilities
        %\item {Advise, review, and test the web site and custom software developed to enhance the management of ASMS membership, conference registration, abstracts, and program.}
\end{tabular}

\mysubsection{Professional Memberships}
\begin{tabular}{V}
  2021--24     & Association of Biomolecular Resources Facilities (ABRF) \\
  2020--     & International Society for Computational Biology (ISCB) \\
  2020--     & Human Proteome Organization (HUPO) \\
  2020--     & United States Human Proteome Organization (US HUPO) \\
  2020--     & American Aging Association (AGE) \\
  2013--     & American Society for Mass Spectrometry (ASMS) \\
\end{tabular}

%\newpage
\mysection{Teaching}

\mysubsection{Teaching Assistance}
\begin{tabular}{V}
    2018 & Introduction to Statistical and Computational Genomics (graduate)
         \newline University of Washington, Seattle \\
    2017 & Fundamentals of Genetics and Genomics (undergraduate)
         \newline University of Washington, Seattle  \\
\end{tabular}

%\newpage

\mysubsection{Quantitative Mass Spectrometry Proteomics Workshop Instructor and Organizer}
\textit{2-5 day workshops held for academic and industry scientists of all backgrounds covering topics including introductory proteomics and mass spectrometry, LC-MS method development, LC-MS data quality control, analysis of quantitative proteomics data} 
\newline \\

\begin{tabular}{V}
  2026 & Claude Code for Scientific Researchers
    \newline virtual \\
  2024 & Ubuntu Proteomics Summer School
    \newline Cradle of Humankind, South Africa \\
  2023--25 & Proteomics (label-free quantification section)
    \newline Cold Spring Harbor Laboratory Courses, New York\\
  2020--25 & Skyline Online
    \newline virtual \\
  2020 & Human Proteome Organization Pre-Congress Training
      \newline virtual \\
  2019 &  Quantitative Targeted Proteomics and Metabolomics Course
      \newline University of Cape Town, Cape Town, South Africa \\
  2017--26 & May Institute on Computation and Statistics for Mass Spectrometry and Proteomics
    \newline Northeastern University, Boston; 2020--21 virtual \\
  2017--19 & Bay-area Targeted Proteomics Course
    \newline Buck Institute for Research on Aging, Novato \\
  2016--25 & Quantitative Targeted Mass Spectrometry Course
    \newline University of Washington, Seattle \\
  2018 &  Skyline@Duke Short Course
      \newline Duke University, North Carolina \\
  2018 &  VIII Proteomics Workshop Skyline
      \newline Center for Research in Energy and Materials, Campinas, Brazil \\
  2017--18 & Case Studies in Quantitative Proteomics Course
    \newline American Society for Mass Spectrometry \\
  2018 & Proteomics 202 Short Course
    \newline Mass Spectrometry: Applications to the Clinical Lab \\

\end{tabular}

%\newpage


\mysection{Audio/Video Features}
\begin{etaremune}

  \item Life Science Washington Institute WIN Mentoring Program (2024) \\ \myref{https://vimeo.com/873842986?fl=pl&fe=sh} 

  \item Genetic Engineering \& Biotechnology News (GEN) LIVE: Is Proteomics the Next Big ``Ome''? (2024) \\ \myref{https://www.genengnews.com/multimedia/is-proteomics-the-next-big-ome/} 

  \item S02E12: THE proteomics BATTLE ROYALE. (2023) The Proteomics Show podcast and US HUPO. \\ \myref{https://open.spotify.com/episode/1rGJGYeTIdqQlCKBqPoUKA?si=0d6eb195498b4f03} \\
  \myref{https://www.youtube.com/watch?v=OybDot4U7uc}
  
  \item Talk is Biotech! with Lindsay Pino. (2022) \\ \myref{https://www.scispot.com/blog/talk-is-biotech-with-lindsay-pino}

  \item Holiday Special - Dr. Lindsay Pino. (2022) The Proteomics Show podcast \\ \myref{https://open.spotify.com/episode/73MEMk6hMGCgfpjIQVqRwh?si=3wni6oXnTAST7zlE-XWV1g}
  
  \item Proteomics calibration with Lindsay Pino. (2021) The Bioinformatics Chat podcast.  \\
    \myref{https://bioinformatics.chat/proteomics-calibration} 
    
  \hide{\item Data independent acquisition mass spectrometry for dynamic spatiotemporal proteomics. (2021) Proteome Software ScaffoldDIA Virtual Brunch featured presentation. \\
    \myref{https://youtu.be/HFDJthoq6hs}}

  \hide{\item Basics of targeted mass spectrometry with Skyline. (2020) Scientific blog post.\\
    \myref{https://lindsaykpino.com/2020/04/29/basics-of-targeted-mass-spectrometry-with-skyline/}}

\end{etaremune}

% Conferences attended without presenting. Kept for the record only.
\hide{\mysection{Conference Attendance}
\begin{tabular}{V}
  2026 & Bioorganic Chemistry Gordon Research Conference \\
\end{tabular}}

%\newpage

\hide{\mysection{References}
{\bf William Noble} \\
Department of Genome Sciences, University of Washington\\
wnoble@uw.edu

{\bf Michael MacCoss} \\
Department of Genome Sciences, University of Washington\\
maccoss@uw.edu

{\bf Susan Abbatiello} \\
Interim Director of the Barnett Institute, Northeastern University\\
su.abbatiello@northeastern.edu}

%\newpage

\end{document}
  